% Options for packages loaded elsewhere
\PassOptionsToPackage{unicode}{hyperref}
\PassOptionsToPackage{hyphens}{url}
\PassOptionsToPackage{dvipsnames,svgnames,x11names}{xcolor}
\documentclass[
  10pt,
  fontset=fandol]{ctexart}
\usepackage{xcolor}
\usepackage[margin=16mm]{geometry}
\usepackage{amsmath,amssymb}
\setcounter{secnumdepth}{-\maxdimen} % remove section numbering
\usepackage{iftex}
\ifPDFTeX
  \usepackage[T1]{fontenc}
  \usepackage[utf8]{inputenc}
  \usepackage{textcomp} % provide euro and other symbols
\else % if luatex or xetex
  \usepackage{unicode-math} % this also loads fontspec
  \defaultfontfeatures{Scale=MatchLowercase}
  \defaultfontfeatures[\rmfamily]{Ligatures=TeX,Scale=1}
\fi
\usepackage{lmodern}
\ifPDFTeX\else
  % xetex/luatex font selection
\fi
% Use upquote if available, for straight quotes in verbatim environments
\IfFileExists{upquote.sty}{\usepackage{upquote}}{}
\IfFileExists{microtype.sty}{% use microtype if available
  \usepackage[]{microtype}
  \UseMicrotypeSet[protrusion]{basicmath} % disable protrusion for tt fonts
}{}
\makeatletter
\@ifundefined{KOMAClassName}{% if non-KOMA class
  \IfFileExists{parskip.sty}{%
    \usepackage{parskip}
  }{% else
    \setlength{\parindent}{0pt}
    \setlength{\parskip}{6pt plus 2pt minus 1pt}}
}{% if KOMA class
  \KOMAoptions{parskip=half}}
\makeatother
\usepackage{longtable,booktabs,array}
\newcounter{none} % for unnumbered tables
\usepackage{calc} % for calculating minipage widths
% Correct order of tables after \paragraph or \subparagraph
\usepackage{etoolbox}
\makeatletter
\patchcmd\longtable{\par}{\if@noskipsec\mbox{}\fi\par}{}{}
\makeatother
% Allow footnotes in longtable head/foot
\IfFileExists{footnotehyper.sty}{\usepackage{footnotehyper}}{\usepackage{footnote}}
\makesavenoteenv{longtable}
\setlength{\emergencystretch}{3em} % prevent overfull lines
\providecommand{\tightlist}{%
  \setlength{\itemsep}{0pt}\setlength{\parskip}{0pt}}
\usepackage{amsmath,amssymb,mathtools,needspace}
\xeCJKDeclareCharClass{CJK}{"2032,"0400->"04FF,"2160->"216B,"2460->"2473,"25A0,"25C7,"25CB,"25CF}
\IfFontExistsTF{Arial Unicode MS}{\xeCJKsetup{AutoFallBack=true}\setCJKfallbackfamilyfont{\CJKrmdefault}{Arial Unicode MS}}{}
\setcounter{secnumdepth}{0}
\setlength{\emergencystretch}{2em}
\allowdisplaybreaks
\usepackage{bookmark}
\IfFileExists{xurl.sty}{\usepackage{xurl}}{} % add URL line breaks if available
\urlstyle{same}
\hypersetup{
  pdftitle={迭代公式的加工},
  colorlinks=true,
  linkcolor={Maroon},
  filecolor={Maroon},
  citecolor={Blue},
  urlcolor={Blue},
  pdfcreator={LaTeX via pandoc}}

\title{迭代公式的加工}
\author{}
\date{}

\begin{document}
\maketitle

\subsubsection{6.2.2 迭代公式的加工}\label{ux8fedux4ee3ux516cux5f0fux7684ux52a0ux5de5}

对于收敛的迭代过程，只要迭代足够多次，就可以使结果达到任意的精度，但有时迭代过程收敛缓慢，从而使计算量变得很大，因此迭代过程的加速是个重要的课题。

\href{../../source-images/pdf-161.jpeg}{原页}

设 \(x_0\) 是根 \(x^*\) 的某个预测值，用迭代公式校正一次得 \(x_1=\varphi(x_0)\)，而由微分中值定理有

\[
x_1-x^*=\varphi'(\xi)(x_0-x^*),
\]

其中 \(\xi\) 介于 \(x^*\) 与 \(x_0\) 之间。

假定 \(\varphi'(x)\) 改变不大，近似地取某个近似值 \(L\)，则由

\[
x_1-x^*\approx L(x_0-x^*),
\]

得

\[
x^*=\frac{1}{1-L}x_1-\frac{L}{1-L}x_0.
\tag{6.2.9}
\]

可以期望，按上式右端求得的

\[
x_2=\frac{1}{1-L}x_1-\frac{L}{1-L}x_0=x_1+\frac{L}{1-L}(x_1-x_0)
\]

是比 \(x_1\) 更好的近似值。

将每得到一次改进值算作一步，并用 \(\bar{x}_k\) 和 \(x_k\) 分别表示第 \(k\) 步的校正值和改进值，则加速迭代计算方案可表述如下：

\[
\begin{aligned}
\text{校正}\quad&\bar{x}_{k+1}=\varphi(x_k),\\
\text{改进}\quad&x_{k+1}=\bar{x}_{k+1}+\frac{L}{1-L}(\bar{x}_{k+1}-x_k).
\end{aligned}
\tag{6.2.10}
\]

\textbf{例 6.5} 求解方程 \(x=\mathrm{e}^{-x}\)。

\textbf{解} 由于在 \(x_0=0.5\) 附近，\((\mathrm{e}^{-x})'\approx-0.6\)，故上述计算公式的具体形式是

\[
\begin{cases}
\bar{x}_{k+1}=\mathrm{e}^{-x_k},\\
x_{k+1}=\bar{x}_{k+1}-\dfrac{0.6}{1.6}(\bar{x}_{k+1}-x_k).
\end{cases}
\]

表 6.5 列出了计算结果。

表 6.5

{\def\LTcaptype{none} % do not increment counter
\begin{longtable}[]{@{}lll@{}}
\toprule\noalign{}
\(k\) & \(\bar{x}_k\) & \(x_k\) \\
\midrule\noalign{}
\endhead
\bottomrule\noalign{}
\endlastfoot
\(0\) & & \(0.5\) \\
\(1\) & \(0.606\,53\) & \(0.566\,58\) \\
\(2\) & \(0.567\,46\) & \(0.567\,13\) \\
\(3\) & \(0.567\,15\) & \(0.567\,14\) \\
\end{longtable}
}

例 6.4 迭代 18 次得到精度 \(10^{-5}\) 的结果 \(x=0.567\,14\)，这里只要迭代 3 次同样得出这个结果。加速的效果是相当显著的。

然而上述加速方案有个缺点，由于其中含有导数 \(\varphi'(x)\) 的有关信息 \(L\)，实际使用不便。

仍设已知 \(x^*\) 的某个猜测值为 \(x_0\)，将校正值 \(x_1=\varphi(x_0)\) 再校正一次，又得

\[
x_2=\varphi(x_1),
\]

由于 \(x_2-x^*\approx L(x_1-x^*)\) 将它与式（6.2.9）联立，消去未知的 \(L\)，有

\[
\frac{x_1-x^*}{x_2-x^*}\approx\frac{x_0-x^*}{x_1-x^*}.
\]

由此推知

\[
x^*\approx\frac{x_0x_2-x_1^2}{x_0-2x_1+x_2}=x_2-\frac{(x_2-x_1)^2}{x_0-2x_1+x_2}.
\]

这样构造出的改进公式确实不再含有关于导数的信息，但是它需要用两次迭代

\begin{quote}
校注（agent 补充）：原页式（6.2.9）确印为等号，转录保留。由于上一步只是用近似常数 \(L\) 得到 \(x_1-x^*\approx L(x_0-x^*)\)，一般只能推出 \(x^*\approx\dfrac{x_1-Lx_0}{1-L}\)；除非该线性误差关系精确成立，式（6.2.9）的等号应作近似关系理解。此为原书等号精度的疑误，非 OCR 误识别。
\end{quote}

\href{../../source-images/pdf-162.jpeg}{原页}

值进行加工。如果将得到一次改进值作为一步，则计算公式如下：

\[
\begin{aligned}
\text{校正}\quad&\tilde{x}_{k+1}=\varphi(x_k),\\
\text{再校正}\quad&\bar{x}_{k+1}=\varphi(\tilde{x}_{k+1}),\\
\text{改进}\quad&x_{k+1}=\bar{x}_{k+1}-\frac{(\bar{x}_{k+1}-\tilde{x}_{k+1})^2}{\bar{x}_{k+1}-2\tilde{x}_{k+1}+x_k}.
\end{aligned}
\tag{6.2.11}
\]

上述处理过程称为 Aitken 方法。

\textbf{例 6.6} 用 Aitken 方法求解方程（6.2.3）。

\textbf{解} 前面曾经指出，求解这一方程的下述迭代公式是发散的：

\[
x_{k+1}=x_k^3-1.
\tag{6.2.12}
\]

现在以这种迭代公式为基础形成 Aitken 算法，即

\[
\tilde{x}_{k+1}=x_k^3-1,\quad
\bar{x}_{k+1}=\tilde{x}_{k+1}^3-1,\quad
x_{k+1}=\bar{x}_{k+1}-\frac{(\bar{x}_{k+1}-\tilde{x}_{k+1})^2}{\bar{x}_{k+1}-2\tilde{x}_{k+1}+x_k}.
\]

仍然取 \(x_0=1.5\)，计算结果如表 6.6 所示。

表 6.6

{\def\LTcaptype{none} % do not increment counter
\begin{longtable}[]{@{}llll@{}}
\toprule\noalign{}
\(k\) & \(\tilde{x}_k\) & \(\bar{x}_k\) & \(x_k\) \\
\midrule\noalign{}
\endhead
\bottomrule\noalign{}
\endlastfoot
\(0\) & & & \(1.5\) \\
\(1\) & \(2.375\,00\) & \(12.396\,5\) & \(1.416\,29\) \\
\(2\) & \(1.840\,92\) & \(5.238\,88\) & \(1.355\,65\) \\
\(3\) & \(1.491\,40\) & \(2.317\,28\) & \(1.328\,95\) \\
\(4\) & \(1.347\,10\) & \(1.444\,35\) & \(1.324\,80\) \\
\(5\) & \(1.325\,18\) & \(1.327\,14\) & \(1.324\,72\) \\
\end{longtable}
}

可以看到，将发散的迭代公式（6.2.12）通过 Aitken 方法处理后，竟获得了相当好的收敛性。在 6.4 节将用几何图像解释这个有趣的现象。

\subsection{本节覆盖原页的校注记录}\label{ux672cux8282ux8986ux76d6ux539fux9875ux7684ux6821ux6ce8ux8bb0ux5f55}

下列记录按原页关联，含同页相邻小节的校注；计算时同时读取上文校注和完整原页。

\begin{itemize}
\tightlist
\item
  \texttt{ch06-E007}：原书 PDF 页 161
\end{itemize}

\end{document}
